\documentclass[11pt]{article}
\usepackage[margin=1in]{geometry}
\usepackage[T1]{fontenc}
\usepackage{lmodern}
\usepackage{amsmath}
\usepackage{xcolor}
\usepackage{listings}
\usepackage{booktabs}
\usepackage{enumitem}
\usepackage{tcolorbox}
\usepackage[hidelinks]{hyperref}

\definecolor{codebg}{RGB}{247,247,249}
\definecolor{kw}{RGB}{0,90,160}
\definecolor{cmt}{RGB}{110,120,110}
\definecolor{str}{RGB}{160,60,60}

\lstset{
  language=Python,
  backgroundcolor=\color{codebg},
  basicstyle=\ttfamily\footnotesize,
  keywordstyle=\color{kw}\bfseries,
  commentstyle=\color{cmt}\itshape,
  stringstyle=\color{str},
  numbers=none, frame=single, framesep=5pt, rulecolor=\color{gray!40},
  breaklines=true, showstringspaces=false, columns=fullflexible,
  xleftmargin=6pt, xrightmargin=6pt, aboveskip=8pt, belowskip=8pt,
}

\newtcolorbox{keybox}[1]{colback=blue!4,colframe=blue!45,title={#1},
  fonttitle=\bfseries,boxrule=0.6pt,arc=2pt}

\title{\vspace{-1.5cm}\textbf{How We Processed the Bonn Dataset}\\[2pt]
\large Its 5 Subsets and How They Became a Binary Task}
\author{EEGdiff\_V2 project}
\date{\today}

\begin{document}
\maketitle
\vspace{-0.6cm}

\section{The Bonn dataset has 5 subsets}
The classic Bonn University EEG database (Andrzejak et al., 2001) ships as
\textbf{5 sets, labeled Z, O, N, F, S} --- 100 single-channel recordings each
(23.6\,s / 4097 samples @ 173.61\,Hz), 500 recordings total. Each set is a
different clinical condition:

\begin{center}
\begin{tabular}{@{}ll@{}}
\toprule
Set (file prefix) & What it actually is \\
\midrule
\textbf{Z} & Healthy volunteers, eyes \textbf{open}, scalp EEG \\
\textbf{O} & Healthy volunteers, eyes \textbf{closed}, scalp EEG \\
\textbf{N} & Epilepsy patients, \textbf{seizure-free interval}, intracranial electrodes in the \\
           & hippocampus \textbf{opposite} to the epileptic focus \\
\textbf{F} & Epilepsy patients, \textbf{seizure-free interval}, intracranial electrodes \textbf{at} the \\
           & epileptic focus \\
\textbf{S} & Epilepsy patients, recorded \textbf{during an actual seizure} (ictal activity) \\
\bottomrule
\end{tabular}
\end{center}

We verified this directly against the raw data: each of the 5 subset zip files
(\texttt{z.zip}, \texttt{o.zip}, \texttt{n.zip}, \texttt{f.zip}, \texttt{s.zip})
extracts to exactly 100 real \texttt{.txt} files (macOS resource-fork junk
excluded), 500 files total.

\section{How we mapped 5 subsets onto our binary task}
Our task is binary seizure-vs-normal, so \texttt{prepare\_bonn()} in
\texttt{dataset\_tools/prepare\_external\_eeg.py} collapses the 5 subsets into
2 classes with one simple rule:

\begin{lstlisting}
label = "seizure" if stem.startswith("S") or parent.startswith("S") else "normal"
\end{lstlisting}

\begin{keybox}{In plain terms}
Only \textbf{Set S} is labeled \texttt{"seizure"}. Sets \textbf{Z, O, N, and F}
are all lumped together as \texttt{"normal"} --- regardless of whether they
are healthy controls (Z, O) or epilepsy patients who simply are not seizing
right now (N, F).
\end{keybox}

We verified the final counts directly from the processed metadata: each of the
5 subsets contributed exactly 100 windows (matching the 100 raw files per
set), and the resulting train/test split is:

\begin{center}
\begin{tabular}{@{}lcc@{}}
\toprule
Split & normal & seizure \\
\midrule
train & 320 & 80 \\
test  & 80  & 20 \\
\midrule
\textbf{total} & \textbf{400} & \textbf{100} \\
\bottomrule
\end{tabular}
\end{center}

That 400/100 = 4:1 ratio is exactly ``4 non-seizure subsets (Z+O+N+F=400) vs.\
1 seizure subset (S=100).''

\section{Worth remembering}
\begin{itemize}[nosep,leftmargin=1.4em]
  \item Our \texttt{"normal"} class is actually a mix of three quite different
  things: healthy eyes-open, healthy eyes-closed, and epilepsy patients who
  are currently seizure-free (sets N and F) --- all treated identically as
  \emph{not seizure}. Clinically, N and F are the \textbf{hardest negatives}
  (epileptic brain tissue, just not actively seizing), so lumping them with
  healthy controls is a real simplification of the task.
  \item This is also why Bonn is a comparatively \textbf{easy} benchmark in
  the literature: the actual discriminative signal (ictal S vs.\ everything
  else) is quite large. Consistent with this, our own result on Bonn was very
  strong (ROC-AUC 0.993).
  \item Each Bonn recording is a \textbf{single channel} with no electrode
  name at all (a separate fact from which subset it belongs to). It is placed
  into cup \#1 (FP1) of our 22-channel frame, with the other 21 channels
  zero-filled --- see the companion handout
  \emph{How We Mapped Every EEG Dataset to 22 Channels}.
\end{itemize}

\end{document}
